\section{Conclusion}\label{sec:conclusion}

This paper studied topology-aware graph contrastive learning for benign-only graph-based intrusion detection. We introduced \sysname{}, a method that learns from benign activity graphs and scores test graphs by distance from benign training embeddings. \sysname{} combines graph-level message-passing representations with topology-aware structural information and performs anomaly scoring in the same contrastive embedding space used during training.

Across three graph-based security datasets, \sysname{} obtains F1 scores of $0.923$ on StreamSpot, $0.790$ on GraSec-IoT, and $0.583$ on Unicorn Wget. It outperforms GraphCL on StreamSpot and Unicorn Wget but underperforms it on GraSec-IoT. The supervised GNN leads on StreamSpot and GraSec-IoT, whereas \sysname{} has the highest F1 on Unicorn Wget. These mixed results show that topology-aware contrastive anomaly scoring can strengthen benign-only detection in some graph settings but is not uniformly beneficial. The result supports continued study of topology-aware detection under incomplete attack labels, while the limited datasets, unmatched baseline configurations, and label-dependent threshold selection motivate broader and more controlled validation.
